\begin{abstract}
Contrary to common intuition, no amount of additional media sampling can recover all four kinetic parameters of a liver-on-a-chip clearance model from media concentration data alone: the limitation is structural, not experimental, and it is independent of any specific platform or drug. Liver-on-a-chip systems are used to estimate intrinsic hepatic clearance (CLint) for in vitro-to-in vivo extrapolation. We analyse a minimal two-compartment model (media and cell, linked by distributional clearance CLd and partition coefficient Kp, with non-specific binding kb and metabolism CLint) under media-only sampling, the most common experimental design. Using the Laplace transform of the system, we prove exactly, in closed form, and independently of sampling density or experiment duration, that only three combinations of the four kinetic parameters are structurally identifiable from media concentration data: denser sampling, longer incubations, or better fitting algorithms cannot close this gap. We give the explicit one-parameter family of parameter quadruples -- parametrized by kb -- that all reproduce an identical media concentration curve, up to numerical solver precision. Independently measuring kb, for example via a cell-free non-specific-binding control already used on some platforms, collapses this family to a single point and renders the remaining parameters structurally identifiable. We then show, by direct computation, that standard local methods for assessing practical identifiability are themselves unreliable near this degeneracy: a Fisher-information covariance approximation has condition numbers of $10^9$--$10^{17}$ and disagrees with a pseudoinverse by up to three orders of magnitude on identical data, and a single-start profile likelihood is contaminated by optimizer non-convergence. Multi-start profile likelihood gives a stable result, and that result is sobering: even with kb and evaporation fixed at their true values, CLint remains weakly constrained by realistic sampling designs. Structural identifiability is therefore necessary but not sufficient for practical estimability, and we propose a minimum reporting standard for liver-chip pharmacokinetic studies that would let readers assess both.
\end{abstract}

\noindent\textbf{Keywords:} liver-on-a-chip; structural identifiability; practical identifiability; microphysiological systems; in vitro-in vivo extrapolation; profile likelihood
